\documentclass[11pt]{article}
\usepackage[utf8]{inputenc}
\usepackage{amsmath}
\usepackage{geometry}
\usepackage{ctex}
\geometry{a4paper, margin=1in}

\begin{document}

\section{现象一：倒置图像的知觉转换}
\textbf{问题：} 为什么某些图片正立时像婴儿面孔，而倒立后却表现为完全不同的景观？

\paragraph{解答：}
这一现象体现了面孔识别的\textbf{整体加工}特性：
\begin{itemize}
    \item \textbf{面孔模板激活：} 当图片正立呈现时，图像中的线条布局符合大脑预设的“面孔图式”，触发了自动化的面孔识别机制，使我们忽略局部细节而感知到整体（婴儿）。
    \item \textbf{结构信息的丧失：} 当图片倒置时，五官之间的空间位置关系不再符合面孔模板。此时大脑被迫采用\textbf{基于特征的加工}，将注意力转向局部的线条和形状，从而识别出原本隐藏在结构中的背景元素（如树木、风景）。
\end{itemize}

\section{现象二：撒切尔效应分析}
\textbf{问题：} 为什么在倒置的面孔上难以察觉局部五官（如眼睛、嘴巴）的反转，而当照片正立时这种扭曲会变得极其惊悚？

\paragraph{解答：}
这进一步证明了人类对面孔的空间配置信息具有极高的敏感度，但这种敏感度具有\textbf{方向依赖性}：
\begin{itemize}
    \item \textbf{倒置状态：} 大脑的面孔模块无法有效运作，识别过程退化为对局部特征的逐一扫描。由于被反转的眼睛和嘴巴在局部特征上仍是“正常的五官”，且相对于倒置的面孔来说，它们的局部朝向（相对于观察者）看似变正了，因此大脑难以察觉其布局的荒谬。
    \item \textbf{正立状态：} 大脑立即启动整体加工，迅速识别出五官之间的\textbf{二阶关系}遭到严重破坏。这种整体结构与预期模板的剧烈冲突，引发了强烈的不适感和恐怖感。
\end{itemize}

\section{结论}
上述两个实验现象共同说明：面孔识别不仅仅是对局部特征（眼、鼻、嘴）的简单加总，更依赖于对五官空间结构的整体把握。而这种高效的整体加工机制主要针对正立的面孔有效。

\end{document}